\begin{question}{Limitations of Transformers on Simple Languages}

\begin{subquestion}{\textbf{Simplifying the attention mechanism.}}

    \begin{equation}
    f' \left( \sum_{t'=1}^{T} s_{t,t'}^{(l)} x_{t'}^{(l)}, x_{t}^{(l)} \right) = f \left( \sum_{t'=1}^{T} s_{t,t'}^{(l)} v_{t'}^{(l)}, x_{t}^{(l)} \right)
    \end{equation}

    \begin{equation}
    f' \left( \sum_{t'=1}^{T} s_{t,t'}^{(l)} x_{t'}^{(l)}, x_{t}^{(l)} \right) = f \left( \sum_{t'=1}^{T} s_{t,t'}^{(l)} ( (W_V^{(l)})^\top x_{t'}^{(l)} ), x_{t}^{(l)} \right)
    \end{equation}

    \begin{equation}
    f' \left( \sum_{t'=1}^{T} s_{t,t'}^{(l)} x_{t'}^{(l)}, x_{t}^{(l)} \right) = f \left( (W_V^{(l)})^\top \left( \sum_{t'=1}^{T} s_{t,t'}^{(l)} x_{t'}^{(l)} \right), x_{t}^{(l)} \right)
    \end{equation}

    \begin{equation}
    f'(z_t^{(l)}, x_t^{(l)}) = f((W_V^{(l)})^\top z_t^{(l)}, x_t^{(l)})
    \end{equation}

\end{subquestion}

\begin{subquestion}{\textbf{Bounding $\|x_{t}^{(l)}\|$.}}

    Proof by induction:

    \subsubsection*{\textbf{Base case}}

    By definition of the row-norm $\|X^{(0)}\|_{2,\infty}$, it holds that:
    \begin{equation}
    \|x_t^{(0)}\| \le \max_{i} \|x_i^{(0)}\| = \|X^{(0)}\|_{2,\infty}
    \end{equation}

    \subsubsection*{\textbf{Inductive step}}

    Assume the bound holds for layer $l$. For layer $l+1$, we use the update rule $x_t^{(l+1)} = f(z_t^{(l)}, x_t^{(l)})$. Since $f$ is Lipschitz continuous with constant $L_f$ and $f(0,0)=0$:
    \begin{equation}
    \|x_t^{(l+1)}\| = \|f(z_t^{(l)}, x_t^{(l)}) - f(0,0)\| \le L_f (\|z_t^{(l)}\| + \|x_t^{(l)}\|)
    \end{equation}

    Because $z_t^{(l)} = \sum_{t'=1}^{T} s_{t,t'}^{(l)} x_{t'}^{(l)}$ and the softmax scores $s_{t,t'}$ are non-negative and sum to 1:
    \begin{equation}
    \|z_t^{(l)}\| = \left\| \sum_{t'=1}^{T} s_{t,t'}^{(l)} x_{t'}^{(l)} \right\| \le \sum_{t'=1}^{T} s_{t,t'}^{(l)} \|x_{t'}^{(l)}\| \le \max_{t'} \|x_{t'}^{(l)}\| = \|X^{(l)}\|_{2,\infty}
    \end{equation}

    Substituting this back into the Lipschitz bound:
    \begin{equation}
    \|x_t^{(l+1)}\| \le L_f (\|X^{(l)}\|_{2,\infty} + \|X^{(l)}\|_{2,\infty}) = 2L_f \|X^{(l)}\|_{2,\infty}
    \end{equation}

    Since this holds for all $t$, we have $\|X^{(l+1)}\|_{2,\infty} \le 2L_f \|X^{(l)}\|_{2,\infty}$. Applying this recursively from $l=0$ yields:
    \begin{equation}
    \|x_t^{(l)}\| \le (2L_f)^l \|X^{(0)}\|_{2,\infty}
    \end{equation}

\end{subquestion}

\begin{subquestion}{\textbf{Bounding Attention Logits.}}

    An entry of the pre-softmax attention scores $l_{t, t'}$ is defined by the dot product of the query $q_t^{(l)}$ and key $k_{t'}^{(l)}$:
    \begin{equation}
    l_{t, t'} = k_{t'}^{(l)\top} q_t^{(l)}
    \end{equation}

    By the Cauchy-Schwarz inequality, the magnitude of this scalar is bounded by the product of the vector norms:
    \begin{equation}
    |l_{t, t'}| \le \|k_{t'}^{(l)}\| \cdot \|q_t^{(l)}\|
    \end{equation}

    Using the definitions $q_t^{(l)} = (W_Q^{(l)})^\top x_t^{(l)}$ and $k_{t'}^{(l)} = (W_K^{(l)})^\top x_{t'}^{(l)}$, we apply the property of matrix operator norms:
    \begin{equation}
    |l_{t, t'}| \le (\|(W_K^{(l)})^\top\| \cdot \|(W_Q^{(l)})^\top\|) \cdot \|x_{t'}^{(l)}\| \cdot \|x_t^{(l)}\|
    \end{equation}

    From sub-question 4b, we have $\|x_t^{(l)}\| \le F$ for all layers $l \in [L]$. Substituting this bound:
    \begin{equation}
    |l_{t, t'}| \le (\|(W_K^{(l)})^\top\| \cdot \|(W_Q^{(l)})^\top\|) \cdot F^2
    \end{equation}

    By defining $L_{att}$ as a constant bounding the product of the query and key weight matrix norms across all layers:
    \begin{equation}
    L_{att} \ge \max_{l} (\|(W_K^{(l)})^\top\| \cdot \|(W_Q^{(l)})^\top\|)
    \end{equation}

    It follows that for any $t, t'$, the magnitude of the logit is bounded by $A$:
    \begin{equation}
    \|l\|_{\infty} \le F^2 L_{att} = A
    \end{equation}

\end{subquestion}

\begin{subquestion}{\textbf{Bounding Softmax Values.}}

    From sub-question 4c, we have $\|l\|_{\infty} \le A$, which implies that for all $j \in \{1, \dots, T\}$:
    \begin{equation}
    -A \le l_j \le A
    \end{equation}
    Applying the monotonically increasing exponential function, we obtain:
    \begin{equation}
    \exp(-A) \le \exp(l_j) \le \exp(A)
    \end{equation}

    The softmax value $s_t$ is defined as the ratio of the exponentiated logit to the sum of all exponentiated logits in the sequence:
    \begin{equation}
    s_t = \frac{\exp(l_t)}{\sum_{j=1}^T \exp(l_j)}
    \end{equation}

    To derive the upper bound for $s_t$, we maximize the numerator ($\exp(A)$) and minimize each term in the denominator's sum ($\exp(-A)$):
    \begin{equation}
    s_t \le \frac{\exp(A)}{\sum_{j=1}^T \exp(-A)} = \frac{\exp(A)}{T \cdot \exp(-A)} = \frac{\exp(2A)}{T}
    \end{equation}

    To derive the lower bound, we minimize the numerator ($\exp(-A)$) and maximize each term in the denominator's sum ($\exp(A)$):
    \begin{equation}
    s_t \ge \frac{\exp(-A)}{\sum_{j=1}^T \exp(A)} = \frac{\exp(-A)}{T \cdot \exp(A)} = \frac{\exp(-2A)}{T}
    \end{equation}

    Combining these inequalities yields the required bound for all $t \in \{1, \dots, T\}$:
    \begin{equation}
    \frac{\exp(-2A)}{T} \le s_t \le \frac{\exp(2A)}{T}
    \end{equation}
\end{subquestion}

\begin{subquestion}{\textbf{Bounding $\Delta_t^{(l)}$.}}

    We define the difference in activations at layer $l$ as $\Delta_t^{(l)} = \|z_t^{(l)} - z_t^{\prime(l)}\|$. 
    Substituting the definition $z_t^{(l)} = \sum_{t'=1}^{T} s_{t,t'}^{(l)} x_{t'}^{(l)}$:
    \begin{equation}
    \Delta_t^{(l)} = \left\| \sum_{t'=1}^{T} s_{t,t'}^{(l)} x_{t'}^{(l)} - \sum_{t'=1}^{T} s_{t,t'}^{\prime(l)} x_{t'}^{\prime(l)} \right\|
    \end{equation}

    Adding and subtracting the term $\sum_{t'=1}^{T} s_{t,t'}^{\prime(l)} x_{t'}^{(l)}$ inside the norm:
    \begin{equation}
    \Delta_t^{(l)} = \left\| \sum_{t'=1}^{T} (s_{t,t'}^{(l)} - s_{t,t'}^{\prime(l)}) x_{t'}^{(l)} + \sum_{t'=1}^{T} s_{t,t'}^{\prime(l)} (x_{t'}^{(l)} - x_{t'}^{\prime(l)}) \right\|
    \end{equation}

    By the triangle inequality and the homogeneity of the norm, we obtain:
    \begin{equation}
    \Delta_t^{(l)} \le \sum_{t'=1}^{T} |s_{t,t'}^{(l)} - s_{t,t'}^{\prime(l)}| \|x_{t'}^{(l)}\| + \sum_{t'=1}^{T} s_{t,t'}^{\prime(l)} \|x_{t'}^{(l)} - x_{t'}^{\prime(l)}\|
    \end{equation}
    which matches equation (49).
\end{subquestion}

\begin{subquestion}{\textbf{Bounding The Difficult, Initial Step.}}

    Using the definition $\|x_{t'}^{(l)}\| \le \|X^{(l)}\|_{2,\infty}$ for all $t'$:
    \begin{align*}
        \sum_{t'=1}^{T} |s_{t,t'}^{(l)} - s_{t,t'}^{\prime(l)}| \|x_{t'}^{(l)}\| &\le \sum_{t'=1}^{T} |s_{t,t'}^{(l)} - s_{t,t'}^{\prime(l)}| \|X^{(l)}\|_{2,\infty} \\
        &\le \left( \sum_{t'=1}^{T} |s_{t,t'}^{(l)} - s_{t,t'}^{\prime(l)}| \right) \|X^{(l)}\|_{2,\infty}
    \end{align*}
\end{subquestion}

\begin{subquestion}{\textbf{Convergence.}}

    Suppose we have two sequences such that $|a_T - b_T| = \mathcal{O}(1/T)$. By definition, there exists a constant $C$ such that for sufficiently large $T$:
    \begin{equation}
    |a_T - b_T| \le \frac{C}{T}
    \end{equation}

    We apply the \textbf{Mean Value Theorem} to the function $f$ on the interval between $a_T$ and $b_T$. This states that there exists some point $c_T$ such that:
    \begin{equation}
    f(a_T) - f(b_T) = f'(c_T)(a_T - b_T)
    \end{equation}

    Taking the absolute value of both sides gives:
    \begin{equation}
    |f(a_T) - f(b_T)| = |f'(c_T)| \cdot |a_T - b_T|
    \end{equation}

    Given that the derivative $f'$ is bounded by some constant $K$ on the relevant domain, we can substitute the known bounds:
    \begin{align*}
    |f(a_T) - f(b_T)| &\le K \cdot |a_T - b_T| \\
    &\le K \cdot \frac{C}{T}
    \end{align*}

    Since $K \cdot C$ is a constant independent of $T$, we conclude:
    \begin{equation}
    |f(a_T) - f(b_T)| = \mathcal{O}\left(\frac{1}{T}\right)
    \end{equation}

    Applying the general result to exp:
    \begin{equation}
    |a_T - b_T| = \mathcal{O}\left(\frac{1}{T}\right) \implies |\exp(a_T) - \exp(b_T)| = \mathcal{O}\left(\frac{1}{T}\right)
    \end{equation}
\end{subquestion}

\begin{subquestion}{\textbf{Softmax.}}

    Skipped.
\end{subquestion}

\begin{subquestion}{\textbf{Bounding The Difficult Part, Final Step.}}

    We assume $\|w_{T,d}^{(l)}\|_{\infty} = \mathcal{O}(1/T)$. This implies the attention logits experience a perturbation of $\mathcal{O}(1/T)$. 

    First, from sub-question \textbf{4b}, the norm of the representations is bounded by a constant $C$ that does not depend on $T$, meaning $\|x_{t'}^{(l)}\| = \mathcal{O}(1)$. 

    Second, applying the result from \textbf{4h}, we know that an $\mathcal{O}(1/T)$ change in logits results in a change to individual softmax weights of $|s_{t,t'}^{(l)} - s_{t,t'}^{\prime(l)}| = \mathcal{O}(1/T^2)$.

    We can now bound the total summation by pulling out the constant term:
    \begin{align*}
    \sum_{t'=1}^{T} |s_{t,t'}^{(l)} - s_{t,t'}^{\prime(l)}| \|x_{t'}^{(l)}\| &\le \left( \max_{t'} \|x_{t'}^{(l)}\| \right) \sum_{t'=1}^{T} |s_{t,t'}^{(l)} - s_{t,t'}^{\prime(l)}| \\
    &= \mathcal{O}(1) \cdot \sum_{t'=1}^{T} \mathcal{O}\left(\frac{1}{T^2}\right) \\
    &= \mathcal{O}(1) \cdot T \cdot \mathcal{O}\left(\frac{1}{T^2}\right) \\
    &= \mathcal{O}\left(\frac{1}{T}\right)
    \end{align*}
\end{subquestion}

\begin{subquestion}{\textbf{Combining All Results}}

We prove Equations (55) and (56) using a single induction on the layer index $l$.

\textbf{Base Case ($l=0$):}

At the input layer, strings $y$ and $y'$ differ only at position $t^*$. Thus, $\|x_{t^*}^{(0)} - x_{t^*}^{\prime(0)}\| = \mathcal{O}(1)$ and $\|x_t^{(0)} - x_t^{\prime(0)}\| = 0$ for $t \neq t^*$. This satisfies the initial condition for (56)[cite: 420, 422].

\textbf{Inductive Step:}

Assume Equation (56) holds for layer $l$. We use this to prove both claims for the next state.

\textbf{Step 1: Proving (55) (Stability of the Attention Output)}
Using the decomposition from 4e [cite: 455-456], we bound the attention output difference $\Delta_t^{(l)}$ as follows:
\begin{equation}
\Delta_t^{(l)} \le \underbrace{\sum_{t'=1}^{T} |s_{t,t'}^{(l)} - s_{t,t'}^{\prime(l)}| \|x_{t'}^{(l)}\|}_{\text{Term 1: Difficult Part}} + \underbrace{\sum_{t'=1}^{T} s_{t,t'}^{\prime(l)} \|x_{t'}^{(l)} - x_{t'}^{\prime(l)}\|}_{\text{Term 2: Residual Errors}}
\end{equation}

\begin{itemize}
    \item \textbf{Term 1 (Application of 4h and 4i):} Sub-question \textbf{4h} establishes that individual attention weights change by $\mathcal{O}(1/T^2)$ [cite: 471-472]. Since representations $\|x_{t'}^{(l)}\|$ are bounded by a constant (from \textbf{4b} [cite: 427-431]), sub-question \textbf{4i} concludes the sum of these $T$ differences is $\mathcal{O}(1/T)$ [cite: 478-479].
    \item \textbf{Term 2 (Application of 4d and Hypothesis):} We split the sum at $t'=t^*$. At the flipped bit, $s_{t,t^*}^{\prime(l)} \|x_{t^*}^{(l)} - x_{t^*}^{\prime(l)}\| = \mathcal{O}(1/T) \cdot \mathcal{O}(1) = \mathcal{O}(1/T)$ using the weight bound from \textbf{4d} [cite: 437-439] and the hypothesis. For $t' \neq t^*$, the hypothesis states $\|x_{t'} - x'_{t'}\| = \mathcal{O}(1/T)$, so the weighted sum is also $\mathcal{O}(1/T)$.
\end{itemize}
Summing these yields $\Delta_t^{(l)} = \mathcal{O}(1/T) + \mathcal{O}(1/T) = \mathcal{O}(1/T)$, proving Equation (55) [cite: 489-491].

\textbf{Step 2: Proving (56) for Layer $l+1$ (Stability of the Representation)}
The next representation is $x_t^{(l+1)} = f(z_t^{(l)}, x_t^{(l)})$[cite: 392]. Applying the Lipschitz continuity of $f$[cite: 398, 412]:
\begin{equation}
\|x_t^{(l+1)} - x_t^{\prime(l+1)}\| \le L_f \left( \Delta_t^{(l)} + \|x_t^{(l)} - x_t^{\prime(l)}\| \right)
\end{equation}
\begin{itemize}
    \item For $t \neq t^*$: Since $\Delta_t^{(l)}$ and the hypothesis $\|x_t^{(l)} - x_t^{\prime(l)}\|$ are both $\mathcal{O}(1/T)$, the result is $\mathcal{O}(1/T)$.
    \item For $t = t^*$: While $\Delta_{t^*}^{(l)} = \mathcal{O}(1/T)$, the term $\|x_{t^*}^{(l)} - x_{t^*}^{\prime(l)}\|$ is $\mathcal{O}(1)$, so the difference remains $\mathcal{O}(1)$.
\end{itemize}
This completes the inductive step, proving that stability is preserved for all layers $l \in [L]$ .

\textbf{Conclusion: Impact on $h(y)$} \\
The final activation is defined as the representation of the EOS token at the final layer, $h(y) = x_T^{(L)}$. Since $t^* < T$ by assumption, position $T$ is not the location of the flipped bit. Thus, applying the bound from Eq. (56) for $t \neq t^*$ at layer $L$, we conclude:
\begin{equation}
\|h(y) - h(y')\| = \|x_T^{(L)} - x_T^{\prime(L)}\| = \mathcal{O}\left(\frac{1}{T}\right)
\end{equation}
\end{subquestion}

\begin{subquestion}{\textbf{Implications on Odd-1.}}

\begin{enumerate}[label=(\roman*), leftmargin=2em, nosep]
\item If $y \in \text{Odd-1}$, flipping a single bit $y_{t^*}$ changes the total count of 1s by exactly one. Consequently, the number of 1s becomes even, meaning $y' \notin \text{Odd-1}$.

\item Because the final activations $h(y)$ and $h(y')$ are nearly identical (differing only by $\mathcal{O}(1/T)$), the resulting probabilities $\hat{p}$ and $\hat{p}'$ will also be nearly identical. The transformer will likely misclassify $y'$ as "Odd" as well.

\item As the string length $T$ approaches infinity, the difference between the representations vanishes ($1/T \to 0$). Thus the model assigns the exact same probability to both the "Odd" and "Even" strings. Since their true labels are different, the transformer will misclassify at least one of them, showing it cannot solve the Odd-1 task for long strings.
\end{enumerate}

\end{subquestion}

\begin{subquestion}{\textbf{A Discussion.}}

\begin{enumerate}[label=(\roman*), leftmargin=2em, nosep]
    \item Layer normalization does not change the fundamental conclusion because it does not fix the "averaging" problem.
    
    \begin{itemize}[leftmargin=1.5em]
        \item \textbf{Case $\epsilon > 0$}: The layer is \textbf{Lipschitz continuous}. Because it is a stable transformation, it cannot amplify a signal that is already vanishing, meaning the difference between representations remains $\mathcal{O}(1/T)$.
        \item \textbf{Case $\epsilon = 0$}: The layer becomes \textbf{numerically unstable}. While this breaks the stability proof, it leads to division by zero or chaotic outputs rather than providing the logic needed to count parity correctly.
    \end{itemize}

    \item The asymptotic result tells us that the model's success on short strings ($|y| \le T^*$) is a \textbf{numerical approximation} rather than a logical understanding of parity. The model cannot learn a scale-invariant rule.

\end{enumerate}

\end{subquestion}
